\documentclass[../../main.tex]{subfiles}
\begin{document}

\begin{center}
  {\fontsize{48pt}{58pt}\selectfont \textbf{2011}\normalsize\textbf{年度}}
\end{center}

\begin{table}[htb]
  \centering
  \caption{出題テーマと難易度}
  \begin{tabular}{|c|p{6.2cm}|c|c|c|}
    \hline
    \textbf{問題番号} & \textbf{テーマ} & \textbf{思考量} & \textbf{計算量} & \textbf{総合} \\
    \hline
    第1問 & 円と直線が作る三角形の面積の最大値 & 易 & 易 & 易 \\
    \hline
    第2問 & 小数部分で定める数列とユークリッドの互除法 & 難 & 普 & 難 \\
    \hline
    第3問 & 円周上を動く点の曲線の長さと極限 & 易 & 普 & 易 \\
    \hline
    第4問 & 放物線上の二点と三角形の重心の軌跡 & 易 & 普 & 普 \\
    \hline
    第5問 & 整数の組の個数の数え上げ & 普 & 普 & 普 \\
    \hline
    第6問 & 二次関数の値域と空間領域の体積 & 普 & 普 & 普 \\
    \hline
  \end{tabular}
\end{table}

\begin{table}[htb]
  \centering
  \caption{解答}
  \renewcommand{\arraystretch}{1.4}
  \begin{tabular}{|c|c|p{9cm}|}
    \hline
    \textbf{問題番号} & \textbf{小問} & \textbf{答え} \\
    \hline
    第1問 & (1) & $S(a)=\dfrac{\sqrt{2a}(1-a)}{a^2+1}$ \\
    \cline{2-3}
    & (2) & $a=2-\sqrt{3}$ \\
    \hline
    第2問 & (1) & $a_n=\sqrt{2}-1$（すべての$n$） \\
    \cline{2-3}
    & (2) & $a=\dfrac{\sqrt{5}-1}{2}$または$\sqrt{2}-1$ \\
    \cline{2-3}
    & (3) & $n\ge q$なら$a_n=0$（証明問題） \\
    \hline
    第3問 & (1) & $u(t)=t\cos\dfrac{L}{t}$，
    $v(t)=t\sin\dfrac{L}{t}$ \\
    \cline{2-3}
    & (2) & $\alpha=\sqrt{L^2+a^2}$，$\beta=\sqrt{L^2+1}$とすると
    $f(a)=\beta-\alpha+L\log\dfrac{L+\alpha}{a(L+\beta)}$ \\
    \cline{2-3}
    & (3) & $-L$ \\
    \hline
    第4問 & -- & $\dfrac{1}{6}<X<\dfrac{1}{2}$，
    $Y=\dfrac{1}{3}\left(\dfrac{2}{6X-1}-\dfrac{1}{4}\right)$ \\
    \hline
    第5問 & (1) & $w=-q$は$p+1$個，$w=p$は$q+1$個 \\
    \cline{2-3}
    & (2) & $0\le s\le p$なら$(p+1)(s+1)$個，
    $p\le s\le2p$なら$(p+1)(2p-s+1)$個，それ以外は$0$個 \\
    \cline{2-3}
    & (3) & $(p+1)^3$個 \\
    \hline
    第6問 & (1) & 二次関数の最大値と最小値の差を場合分けして求める（解答参照） \\
    \cline{2-3}
    & (2) & 領域$S$を図示（解答参照） \\
    \cline{2-3}
    & (3) & $\dfrac{17}{18}$ \\
    \hline
  \end{tabular}
\end{table}

\paragraph{試験形態}
6問．試験時間は要確認．

\end{document}
